\section{Factorization of the finite sieve coefficients}
\label{app:coefficients}

We prove that the actual upper and lower sieve coefficients have
the factorization required by the distribution estimate. Throughout
this appendix $N$ is the even target. Finite partitions, levels,
and approximation tolerances are chosen before $N$ tends to
infinity; the arithmetic coefficients themselves may depend on $N$.
We write $\omega_{\mathrm{pf}}(n)$ for the number of distinct prime
factors of $n$.

\subsection{Finite policies and residual square conditions}

Write
\[
 \rho_0=\frac{19}{96},\qquad \mu_0=\frac{25}{128},\qquad
 \Lambda_{\rm cap}=\frac{153}{256},\qquad \Theta=\frac{299}{500}.
\]
The number $\Theta$ bounds the complete carrier; $\Lambda_{\rm cap}$ bounds the
high-rank levels.  A finite policy specifies a level $N^\lambda$ at each
of finitely many prime prefixes.  An ancestor decision is fixed before
the children are subdivided.  For decreasing exponents $t_1,\ldots,t_j$,
put $S_j=t_1+\cdots+t_j$.  The full square conditions at an auxiliary
level $N^\ell$ are
\begin{equation}
 S_{j-1}+2t_j\le\ell.
 \label{coef:full-squares}
\end{equation}
If $I$ is a selected subset of the marked slots, put
$b_I=\sum_{i\in I}t_i$.  In addition to \eqref{coef:full-squares}, retain
all the conditions
\begin{equation}
 b_I+\sum_{\substack{i<j\\i\notin I}}t_i+2t_j\le\ell
 \qquad(j\notin I).
 \label{coef:selected-squares}
\end{equation}
For $j$ after the last selected index, the left side of
\eqref{coef:selected-squares} is exactly $S_{j-1}+2t_j$.
Thus adjoining smaller primes preserves the selected conditions whenever
the full conditions are preserved.  An admitted cube also implies the
next square, since
\[
 S_j+2t_{j+1}\le S_{j-1}+3t_j.
\]
This observation applies to both parities of the ordinary sieve tree:
upper first exits have odd rank, and lower first exits have even rank.

The physical level $\lambda$ and the auxiliary covering level $\ell$
have different uses.  The coefficient and its density use $\lambda$;
factorization may use a slightly larger $\ell$.  No change in $\lambda$
is made in obtaining this cover.  Here is the needed choice at the
three-slot alternative.  Put
\[
 A(t)=\frac{1+t}{2},\qquad W(z)=\Lambda_{\rm cap}-\frac38z.
\]
On a coarse box let $\Gamma=\inf\max(A,W)$ and
$\lambda=\Gamma-2e$, where the variation allowance satisfies
$\chi<e/10$.  If $\inf A\ge\lambda+e/2$, choose
$\ell=\lambda+e/4$.  Otherwise $\sup A<\lambda+11e/20$,
so $W\ge\Gamma$ throughout the box; choose $\ell=\lambda+e$.
This leaves a reserve of at least $9e/20$ above the lower bound and $e$
below the upper bound.  The remaining initialization reserves include
\[
 1-\Lambda_{\rm cap}-2\mu_0=\frac3{256},\qquad
 \frac{525}{8704}+\frac{49z}{272}>0.
\]
The higher-rank changes are restricted to compact pieces on which their
support conditions have positive reserve.  Fine endpoint enlargement is then taken
smaller than this reserve.

At an ordinary lower start $t_1\le {h_{\rm s}}=19/125<\mu_0$, the endpoint
admissibility condition satisfies
\[
 \max\left(\frac{1+t_1}{2},\Lambda_{\rm cap}-\frac38t_1\right)
 \ge\frac{249}{448}.
\]
A level with a fixed buffer below $249/448$ therefore has a strict
covering interval.  A restored fourth exit is admitted only after its
own cube and all the full and selected squares have been checked.
For upper prefixes above $\mu_0$, the function
\[
 v(x)=\min\left(\frac{1+x}{2},\frac{57}{64}-\frac32x\right)
\]
satisfies $v(x)>3x$ for $x<\rho_0$.  Selecting $y<x$ gives
$2x+y<3x$; selecting $y+z<x$ gives $2x+y+z<3x$.
These are the initial square bounds for the two additional alternatives.

In a marked state write $S=\log_N B$, $L=\lambda-S$ and
$\alpha=(1-S)/L\ge2$.  A singleton of exponent $Lu$ changes the
residual length to $L(1-u)$ and a pair of exponents $Lu,Lv$ changes it
to $L(1-u-v)$.  In both cases $\alpha$ increases.  The conditions
$2u<1$, or $2v<1$ and $v+2u<1$, are precisely the new full square
conditions.  The selected conditions then follow from
\eqref{coef:selected-squares}; the parent prime ceiling is retained.
The finite construction has at most $440$ elementary transitions.
Its residual exponents are at least $L_{\min}=\ell_0 101^{-440}$
for some $\ell_0>0$, and its appended prime exponents are at least
$L_{\min}/6$.  In particular all switched leaves have a common positive
prime cutoff.

\subsection{Finite approximation and strict margins}

The finite partition is chosen jointly for all the marked queries.
At the third upper exit, the omitted last-prime tail is
\[
 O\!\left(\eta\left(1+\log^2\frac{\rho_0}{\eta}\right)\right),
 \qquad\eta\downarrow0.
\]
At lower exits of rank $r=3,4,5$, the corresponding bound is
$O(\eta(1+\log^{r-1}({h_{\rm s}}/\eta)))$.  The higher changes have positive
last-prime cutoffs; in particular the eighth-rank domains have
$t_8>53/13824$.  On the compact remainder, proper affine walls can be
trimmed with arbitrarily small mass.  A region on which the chosen
label is identically inherited is retained rather than treated as a
wall.  Freezing the coarse decision before refining the children gives
one common finite policy for the count and the density.

For the third envelope, the coarse parameters $e,\chi$ give
$0\le L_{\mathrm{end}}-\lambda\le4e+6\chi$.
Fix the finite policy and its positive reserves. The admissibility
functions are continuous on finitely many compact boxes, so their
strict reserves have a positive minimum. Hence some $\kappa_*>0$
preserves all the inequalities when
$\log_xN\in[1,1+\kappa_*]$ and the endpoint enlargement is
sufficiently small. The margin is obtained on these fixed compact
pieces, after the limiting boundary regions have been removed.

\subsection{Three-factor representations}

We describe the endpoint construction, since the factorization must hold
for the coefficient as a whole.  Let $x$ be the complete local scale,
$q_0=\log_x N$, $\vartheta=7/32$, and $\alpha_0\ge q_0$ the residue
parameter in the local admissibility condition.  For a small fixed $\zeta>0$ put
\[
 \theta_{\rm cov}=q_0\ell+2\zeta,\qquad D=x^{\theta_{\rm cov}-2\zeta},
 \qquad A=Ux^{-\zeta},
\]
where $U$ is the chosen medium length.  Define
\begin{align*}
 v_0&=2\theta_{\rm cov}-1,\\
 u_0&=\min\left\{
 \frac{2-(2+2\vartheta)\theta_{\rm cov}}{3-4\vartheta},
 \frac{2-2\theta_{\rm cov}-\alpha_0\vartheta}{3}\right\}.
\end{align*}
In this subsection the slot exponents and $b_I$ are measured to base
$x$: they are $q_0$ times their preceding base-$N$ values.  The covering
exponent in the full and selected square conditions is therefore
$q_0\ell$.
The selected subset admissibility condition is $v_0\le b_I\le u_0$,
with \eqref{coef:full-squares}--\eqref{coef:selected-squares} imposed
on ordered upper endpoints.  All lines of this admissibility condition are retained
after rescaling, including a line that is dominated when $\alpha_0=1$.

Suppose first that $C=D/(Ax^{b_I})\ge1$.  Put the selected slots in
group two and insert the remaining slots greedily into capacities $A,C$.
At every insertion the selected square bounds the square of the new
slot by the product of the remaining capacities, so at least one accepts
it.  The resulting group endpoints satisfy
\[
 B_1\le A,\qquad B_2=x^{b_I},\qquad B_3\le C
\]
and
\begin{align}
 A^2B_2B_3^2&\le x^{1-4\zeta},\notag\\
 N^\vartheta AB_1B_2^5B_3^2&\le x^{2-4\zeta},
 \label{coef:three-monomials}\\
 A^{1+2\vartheta}B_1B_2^{5-2\vartheta}B_3^{2+2\vartheta}
 &\le x^{2-(4+4\vartheta)\zeta}.\notag
\end{align}
If $C<1$, insert all slots into $A,D/A$ using the full squares.
Here $1<D/A<x^{b_I}$.  Set $B_3=1$ and replace $x^{b_I}$ by $D/A$
in the last two bounds.  The first bound follows from
$A^2B_2\le AD\le x^{23/24-5\zeta/2}$.  Thus no factor with an
endpoint below one is introduced.

For the three-slot $W$ alternative initialize capacities
$A,x^{v_0},D/(Ax^{v_0})$ with the first two slots.  One fits $x^{v_0}$;
the square of the other is at most the product
$x^{1-\theta_{\rm cov}-2\zeta}$ of the two remaining capacities.  Continue
inserting whole slots.  If a first slot fits none of the capacities,
adjoin it to group two.  Its new product $B$ lies in
$(x^{v_0},x^{u_0}]$, since that slot is at most the third slot and
$p_3\le x^{u_0-v_0}$.  Failure of the first capacity together with the
full square gives $ABc<D$, so the two residual capacities are legal.
The full-square greedy argument finishes the partition.  If no failure
occurs, the capacity product gives \eqref{coef:three-monomials} directly.
These algorithms use only endpoint data of a whole atom and its medium
length.  They never use the individual arithmetic tuple or the sign of
a progression discrepancy.

\subsection{The coefficient identity and squarefree restriction}

Refine the interval labels to a common disjoint half-open alphabet.
If $m_I$ slots have label $I$, set $M=\prod_I m_I!$.
The unrestricted prime convolution $N_J$ of the slot list has value
$M/\prod_p\nu_p(n)!$ on its support.  Its sorted squarefree coefficient
is therefore $\one_{\mathrm{sf}}N_J/M$.  For a fixed three-group
partition put $M_j=\prod_I m_{I,j}!$ and let $\gamma_j$ be the group
convolution divided by $M_j$.  Then $0\le\gamma_j\le1$ on every integer
and the exact identity is
\begin{equation}
 \frac{N_J}{M}
 =\frac{M_1M_2M_3}{M}\,
       (\gamma_1*\gamma_2*\gamma_3).
 \label{coef:tensor-identity}
\end{equation}
An empty group is the unit mass at one.  A marked convolution already
contains its marked multiplicities.  Only its ordered residual slots
are divided by their factorials: regrouping a signed atom
$\epsilon\one_{\mathrm{sf}}N_{\mathrm{all}}/M_{\mathrm{res}}$
gives the external scalar $\epsilon M_1M_2M_3/M_{\mathrm{res}}$.
This scalar is fixed, although it need not be bounded by one.

Choose $w=N^{\eta_0}$ below every marked prime exponent.  For $p<w$,
$p\mid BN$ if and only if $p\mid N$. Hence one pair of small-prime
beta weights, with support below the smallest residual reserve, serves
every state in the finite marked family.  The sign of the atom determines
which member is attached.  The small and rough prime universes are
disjoint, so decomposition into these two parts is unique.  Attaching
the small-prime beta weight to one of the groups in
\eqref{coef:tensor-identity} preserves one-boundedness.

Global squarefreeness remains in the physical coefficient and its density.
It is removed only in the analytic tensor expansion.  If the support is
$D_*\le N^{\theta_*}$ with $\theta_*<1$, the change in the discrepancy
is at most a fixed divisor factor times
\begin{align}
 &\sum_{\substack{q\le D_*\\p^2\mid q\text{ for some }p\ge N^{\eta_0}}}
 \left(\frac{N}{q}+1+\frac{N}{\varphi(q)}\right)
 \ll (N+D_*)N^{-\eta_0+o(1)}.
 \label{coef:collision}
\end{align}
The term $1$ includes an exact complete product equal to $N$.
No divisor function is evaluated at zero.

\subsection{Parameter order and the medium estimate}

Before small-prime attachment, endpoint enlargement and local rescaling,
the complete supports at rough rank at most two have exponents
\[
 \frac{19}{48},\qquad \frac{38}{125},\qquad
 \frac{1527}{3584}
\]
for the upper, ordinary lower and marked families, respectively.
Their gaps below $1/2$ are $5/48$, $49/250$ and $265/3584$.
For the last bound use $x<57/224$, $x+2y\le\Lambda_{\rm cap}$, and hence
$x+y\le(\Lambda_{\rm cap}+57/224)/2=1527/3584$.
The fixed coefficient norm permits the below-half-level
Bombieri--Vinogradov estimate for these pieces.  A high-rank coefficient
requires the factorization above.

First fix the finite collection of estimates and their dependencies, coarse policy, factorization conditions and
positive reserves.  Choose $\zeta$ below these reserves and the prime
cutoffs, and set $\sigma=\zeta/2000$.  Next choose a common beta support
small enough for the density and support tolerances.  If the rough-rank
bound is $J$, take the endpoint mesh $\tau$ so that $(J+2)\tau$ is below
every remaining reserve.  Only after the fixed divisor and moment orders
are known choose
\begin{equation}
 0<\kappa\le\min\left\{\kappa_*,\frac{\sigma}{200},
                         \frac1{2K_{\mathrm{mom}}}\right\},
 \qquad
 0<\kappa_0<\frac{\kappa}{4(1+\kappa)}.
 \label{coef:parameter-order}
\end{equation}
Here $K_{\mathrm{mom}}$ bounds the fixed moment orders in the Fourier
tail and complementary-coefficient truncation.  Make $\kappa_0$ smaller
if necessary so that the total carrier, safely bounded by
$\theta_*=599/1000$, satisfies $\theta_* /(1-\kappa_0)<3/5$ and all
three low-rank supports remain below one half.

Replacing $A$ by $U$ in \eqref{coef:three-monomials} leaves savings
$\zeta,2\zeta,9\zeta/16$.  After two reserves of $\zeta/16$, the
smallest saving is $7\zeta/16>15\sigma$.  Small-prime attachment enlarges
only group three; its costs $2\ell_\beta,2\ell_\beta,
(2+2\vartheta)\ell_\beta$ are included in the reserves.  Dyadic subdivision
therefore produces $Q_1,Q_2,Q_3$ satisfying all four size conditions of
Proposition~\ref{prop:current-medium}.  The strict margins absorb fixed
endpoint constants.  There are finitely many atoms, with fixed total
absolute scalar norm, before $N$ grows.  This establishes a uniform
three-factor representation for the actual finite coefficient families.

\hypertarget{coef:medium-checks}{}
\paragraph{Checks for the medium branch.}
The following checks connect Proposition~\ref{prop:current-medium} with
its application in Proposition~\ref{coef:prime-product-insertion}. They are
made after fixing the common finite policy and its strict support reserves.
\begingroup
\renewcommand{\LTcaptype}{}
\stepcounter{LT@tables}
\small
\setlength{\tabcolsep}{4pt}
\begin{longtable}{@{}>{\raggedright\arraybackslash}p{.21\textwidth}>{\raggedright\arraybackslash}p{.45\textwidth}>{\raggedright\arraybackslash}p{.28\textwidth}@{}}
\toprule
 Condition or input & Realization in the application & Verification \\
\midrule
\endhead
\bottomrule
\endfoot
Complete scale and residue & $x=UV$ is the scale of the complete product and $\mathfrak a=N$. The trim and local rescaling give $c_0x\le N\le x^{1+\kappa}$. & \eqref{coef:complete-scales}, \eqref{coef:parameter-order}\\[3pt]
Rough medium sequence & An individual interval carries the rough $1$ or M\"obius sequence, its bounded weight and permitted Mellin twist; the auxiliary-coprimality estimate is retained. & \eqref{coef:rough-SW} and its proof\\[3pt]
Complementary sequence & Unique factorization gives a fixed divisor bound. Any logarithmic truncation is charged by the mixed moment estimate before applying the local bound. & Lemma~\ref{coef:mixed-moment}; proof of Proposition~\ref{coef:prime-product-insertion}\\[3pt]
Three modulus factors & The normalized group convolutions are separate and one-bounded, with the residue-coprimality support retained. Marked multiplicities remain external scalars. & \eqref{coef:tensor-identity}; full and selected squares \eqref{coef:full-squares}--\eqref{coef:selected-squares}\\[3pt]
Four strict size tests & The three monomial bounds, the medium capacity, small-prime attachment and local rescaling yield all the tests in \eqref{eq:current-medium-ranges}. & \eqref{coef:three-monomials}; the reserves following \eqref{coef:parameter-order}\\[3pt]
Moving-residue phase and spectral input & The exact reciprocal phase is retained. Its variation and the packet and spectral norms use the stated local margins. & \eqref{eq:current-retained-phase}, \eqref{ana:phase-variation}; Appendix~\ref{app:classical}\\[3pt]
Uniform summation errors & Fixed atom scalars, Fourier norms and logarithmically many charts are covered by stronger local logarithmic accuracy. Prime powers, squarefree collisions and exact products equal to $N$ are counted separately. & \eqref{coef:collision}; proof of Proposition~\ref{coef:prime-product-insertion}\\[3pt]
\end{longtable}
\endgroup

For the four size tests, deducting the stated support reserves from the
savings $\zeta,2\zeta,9\zeta/16$ leaves strict margins that absorb fixed
dyadic endpoint multiples. This check concerns the medium branch; the
below-half estimate and Appendix~\ref{app:brun} cover the other branches of
Proposition~\ref{coef:prime-product-insertion}. Upper coefficients, lower
coefficients and ordinary marked descendants enter through that proposition
with their actual signed coefficient $\rho_N$.

\section{Distribution of the smooth prime products}
\label{app:prime-products}

This appendix proves \eqref{eq:current-PP} for the finite coefficient families of Appendix~\ref{app:coefficients}. The medium branches use Proposition~\ref{prop:current-medium}; the remaining branches use the below-half estimate or Appendix~\ref{app:brun}.

For a smooth weight $W_N$ on at most $k_*$ labelled prime slots, with
$p_i\ge cN^{a_*}$ and product $O(N)$, define
\[
 E_q(N;W)=\sum_{p_1,\ldots,p_k}W_N(p_1,\ldots,p_k)
 \left(\one_{p_1\cdots p_k\equiv N\ (q)}
       -\frac{\one_{(p_1\cdots p_k,q)=1}}{\varphi(q)}\right).
\]
The geometric weights have fixed polylogarithmic mixed Euler derivative
bounds.  They carry no additional arithmetic filters.

\subsection{Two auxiliary mean estimates}

Let $d_j$ denote the $j$-fold divisor function.
We record the elementary moment estimate with its dependence on the
moving shift, since logarithmic truncations require a fixed moment order.

\begin{lemma}\label{coef:mixed-moment}
Fix nonnegative integers $R,S$ and a constant $C\ge1$.  Put
$K_{\rm mom}=2^{2\max(R,S,1)}$ and $0\le\kappa\le1/(2K_{\rm mom})$.
Uniformly for $|a|\le Cx^{1+\kappa}$ and an interval
$J\subseteq[1,Cx]$,
\[
 \sum_{\substack{b\in J\\b\ne a}}
       \tau(b)^R\tau(|b-a|)^S
 \ll_{C,R,S} x(1+\log x)^{K_{\rm mom}-1}.
\]
\end{lemma}

\begin{proof}
If an interval $I\subseteq[1,Z]$ has length $Y$, choosing a largest
coordinate of an ordered $j$-tuple whose product belongs to $I$ gives
\begin{align*}
 \sum_{n\in I}d_j(n)
 &\le j\sum_{d\le Z^{1-1/j}}d_{j-1}(d)(Y/d+1)\\
 &\le j\{Y(1+\log Z)^{j-1}
              +Z^{1-1/j}(1+\log Z)^{j-2}\}.
\end{align*}
Indeed the product of the other $j-1$ coordinates is at most
$Z^{1-1/j}$; dropping the largest-coordinate restriction only increases
the sum.  The two final bounds follow by expanding the divisor function
as an ordered product and summing the last factor.
For every integer $e\ge0$,
$\tau(n)^e\le d_{2^e}(n)$.  On a prime power this follows by mapping
an $e$-tuple in $\{0,\ldots,v\}^e$ to the histogram of the subsets
of coordinates reaching each of the positions $1,\ldots,v$; the
histogram is a weak composition of $v$ into $2^e$ parts and determines
the tuple.

The nonzero values $|b-a|$ lie, with multiplicity at most two, in at most
two intervals of length $O_C(x)$ and upper endpoint $O_C(x^{1+\kappa})$.
Apply the preceding interval bound to $d_{K_{\rm mom}}$ to obtain
\[
 \sum_{\substack{b\in J\\b\ne a}}\tau(|b-a|)^{2S}
 \ll x(1+\log x)^{K_{\rm mom}-1}
       +x^{(1+\kappa)(1-1/K_{\rm mom})}(1+\log x)^{K_{\rm mom}-2}.
\]
The second exponent is at most $1-1/(2K_{\rm mom})-1/(2K_{\rm mom}^2)<1$.
The unshifted $2R$ moment has the same bound.
Cauchy--Schwarz proves the assertion.  The term $b=a$ is removed
before applying the divisor function.
\end{proof}

\begin{lemma}[A convolution estimate below one half]
\label{coef:low-rank-BV}
Fix $\eta>0$ and $c\ge1$. Let $x=UV$,
$x^\eta\le U,V\le x^{1-\eta}$, and let $\alpha,\beta$ be supported
on $[U,cU]$ and $[V,cV]$ with
\[
 \sum_m|\alpha(m)|^2\ll U(\log x)^C,\qquad
 \sum_n|\beta(n)|^2\ll V(\log x)^C.
\]
They are independent of the summation modulus.  Suppose that for every
fixed $A_1,B>0$, every $k\ge1$, $r\le(\log x)^B$ and $(b,r)=1$,
\begin{equation}
 \left|\sum_{\substack{m\equiv b\ (r)\\(m,k)=1}}\alpha(m)
 -\frac1{\varphi(r)}\sum_{(m,rk)=1}\alpha(m)\right|
 \ll_{A_1,B}\tau(k)^{C_0}U(\log x)^{-A_1},
 \label{coef:small-aux-SW}
\end{equation}
with one fixed exponent $C_0$.  For every $A>0$ there is a $D>0$ such
that $Q\le\sqrt x(\log x)^{-D}$ implies
\[
 \sum_{q\le Q}\max_{(a,q)=1}
 \left|\sum_{mn\equiv a\ (q)}\alpha(m)\beta(n)
 -\frac1{\varphi(q)}\sum_{(mn,q)=1}\alpha(m)\beta(n)\right|
 \ll_A x(\log x)^{-A}.
\]
\end{lemma}

\begin{proof}
Put $L=\log x$.  The primitive multiplicative large sieve
\cite[Theorem~18.2]{Kedlaya} is
\begin{equation}
 \sum_{r\le R}\frac r{\varphi(r)}\sum_{\chi\ (r)}^{*}
 \left|\sum_nz(n)\chi(n)\right|^2
 \le(H-1+R^2)\sum_n|z(n)|^2,
 \label{coef:multiplicative-LS}
\end{equation}
when $z$ is supported on an interval of at most $H$ consecutive integers.
Character orthogonality removes the principal character exactly.
A character modulo $q=rk$ induced by a primitive character $\chi^*$
satisfies $\chi(m)=\chi^*(m)\one_{(m,k)=1}$, including when $(r,k)>1$.
Put $\alpha_k=\alpha\one_{(m,k)=1}$ and
$\beta_k=\beta\one_{(n,k)=1}$ and write $A_k(\chi),B_k(\chi)$ for
their character sums.  Since $\varphi(rk)\ge\varphi(r)\varphi(k)$,
the required sum is bounded by
\begin{equation}
 \sum_{k\le Q}\frac1{\varphi(k)}
 \sum_{1<r\le Q/k}\frac1{\varphi(r)}
              \sum_{\chi\ (r)}^{*}|A_k(\chi)B_k(\chi)|.
 \label{coef:induced-characters}
\end{equation}
Both coprimality restrictions are essential.  The restrictions do not
increase the $L^2$ norms and, for fixed $k$, do not depend on $r$.

Let $R_0=L^{B_0}$.  For $R<r\le2R$, $R\ge R_0$, Cauchy--Schwarz
and \eqref{coef:multiplicative-LS} bound the inner sum by
\[
 \ll L^{C'}\left\{\frac xR+
 \sqrt x(\sqrt U+\sqrt V)+R\sqrt x\right\}.
\]
Summing dyadically, then over $k$ in
\eqref{coef:induced-characters}, gives
\begin{equation}
 \ll L^{C'}\left\{\frac{xL}{R_0}
       +L^2\sqrt x(\sqrt U+\sqrt V)+Q\sqrt x\right\}.
 \label{coef:large-conductors}
\end{equation}
Here $\sum_{k\le Y}1/\varphi(k)\ll1+\log Y$ and
$\sum_{k\ge1}1/(k\varphi(k))<\infty$; the latter sum occurs in
the last term.  Both follow from
$n/\varphi(n)=\sum_{d\mid n}\mu(d)^2/\varphi(d)$ and
$1/(d\varphi(d))\le\tau(d)/d^2$.
The middle term in \eqref{coef:large-conductors} has a fixed power
saving because $\sqrt x(\sqrt U+\sqrt V)\le2x^{1-\eta/2}$.
Taking $B_0,D$ sufficiently large controls the other two terms.

For $r\le R_0$, sum \eqref{coef:small-aux-SW} against the primitive
nonprincipal character.  Its mean vanishes, yielding
$|A_k(\chi)|\ll\varphi(r)\tau(k)^{C_0}UL^{-A_1}$.
The $L^1$ bound for $\beta_k$ is $O(VL^{C/2})$, so the small-conductor
part of \eqref{coef:induced-characters} is at most
\[
 \ll xL^{C/2-A_1}
       \left(\sum_{r\le R_0}\varphi(r)\right)
       \left(\sum_{k\le Q}\frac{\tau(k)^{C_0}}{\varphi(k)}\right).
\]
Choose an integer $J\ge C_0+1$.  Since
$n/\varphi(n)\le\tau(n)$ and $\tau(n)^J\le d_{2^J}(n)$,
the last sum is at most $(1+\log Q)^{2^J}$.
After $B_0$ is fixed, taking $A_1$ sufficiently large finishes the proof.
All divisor orders are fixed before the requested logarithmic accuracy.
\end{proof}

For an original prime slot, the classical prime Siegel--Walfisz theorem
\cite[Exercise~64]{TaoSW}, followed by partial summation, gives
\eqref{coef:small-aux-SW} at $k=1$ with any prescribed saving.
For arbitrary $k$ the change deletes only primes dividing $k$ and is at
most $2\|\alpha\|_\infty\omega_{\mathrm{pf}}(k)$; its fixed logarithmic norm is
absorbed by the positive-power scale $U$.  Thus $C_0=1$ works for all
the required exponents.  Fixed smooth profiles, interval endpoints and
Mellin twists with fixed polylogarithmic variation preserve the estimate.
For several original primes, put one such slot in $\alpha$ and the
remaining factors in $\beta$.  Unique factorization supplies a fixed
divisor bound and hence the required $L^2$ norm for $\beta$.

\subsection{Proof of the prime-product estimate}

\begin{proposition}\label{coef:prime-product-insertion}
For the finite upper, lower and marked coefficients constructed in
Appendix~\ref{app:coefficients}, the restricted medium estimate of
Proposition~\ref{prop:current-medium} and the below-half-level
Bombieri--Vinogradov estimates imply
\[
 \sum_q\rho_N(q)E_q(N;W)=O_{A,W}\bigl(N(\log N)^{-A}\bigr)
 \qquad(A>0).
\]
The coefficient in the count and exact coprime mean is the same.
The estimate is uniform for the fixed finite family and all sufficiently
large even $N$.
\end{proposition}

\begin{proof}
Choose a smooth $\eta$ equal to zero on $[0,1]$ and one on $[2,\infty)$,
and multiply $W_N$ by $\eta(p_1\cdots p_k/N^{1-\kappa_0})$.
The discarded products have size less than $2N^{1-\kappa_0}$, and are
strictly below $N$ for large $N$.  Their complete modulus-divisor sum
and exact mean have error $O(N^{1-\kappa_0+o(1)})$.
All Euler derivatives of this trim are bounded.

Insert a smooth nonnegative dyadic partition in each original prime.
On an original chart put $R_i=2^{j_i}$, $X=\prod_iR_i$ and $X^+=2^kX$.
With $f(v)=v^{1/(\log\log v)^3}$ fix the common roughness threshold
\[
 z=f(2^{6k_*+2}X^+)
\]
before applying a convolution identity.  Every original prime exceeds
$z$ eventually.  Replacing its prime indicator by $\Lambda(n)/\log n$
and deleting proper prime powers costs a fixed power saving: their
reciprocal mass above $cN^{a_*}$ is $N^{-a_*/2+o(1)}$.
Nonzero complete differences have divisor cost $N^{o(1)}$, and an exact
product $N$ contributes at most $N^{\theta_*+o(1)}$ separately.

For a group at scale $R$ let $Z=(2R)^{1/3}$ and
$\mu_Z=\mu\one_{[1,Z]}$.  If $\delta$ denotes the convolution identity,
$\delta-1*\mu_Z$ is supported above $Z$; consequently
$\Lambda*(\delta-1*\mu_Z)^{*3}=0$ below $2R$.  Expanding gives
\begin{equation}
 \Lambda=\sum_{j=1}^3(-1)^{j-1}\binom3j
             \mu_Z^{*j}*1^{*(j-1)}*\log.
 \label{coef:HB}
\end{equation}
The original smooth chart is retained.  The roughness projector factors
over all the variables in \eqref{coef:HB}; use $P^-(1)=\infty$.
There is no additional sharp coupled group boundary.

Each original group has at most six variables.  Individual rectangles
meeting its smooth support give complete local scales satisfying
\begin{equation}
 2^{-7k-1}X\le x\le2^{6k-1}X^+,
 \qquad f(x)\le z\le x^{2/(\log\log x)^3}.
 \label{coef:complete-scales}
\end{equation}
Thus $cN^{1-\kappa_0}\le x\le CN$ and
$c_0x\le N\le x^{1+\kappa}$ by \eqref{coef:parameter-order}.
The scale $x$ is always that of the complete product.

To separate the geometric profile, use coordinates
$u_j=\log(n_j/U_j)$ and a smooth compact extension $w$ of the profile.
Writing
$D_w=\sum_{S\subseteq\{1,\ldots,m\}}\|\partial_{2S}w\|_1$,
integration by parts gives
\[
 |\widehat w(\xi)|\le D_w\prod_j(1+\xi_j^2)^{-1},
 \quad \frac1{(2\pi)^m}\int|\widehat w(\xi)|\,d\xi
 \le\frac{D_w}{2^m}.
\]
The tail outside $|\xi|_\infty\le H$ is $O(mD_w/H)$.
Here $D_w$ is a fixed logarithmic power.  The arithmetic sums can be
interchanged with the absolutely integrable Fourier representation.
After the complete modulus sum, the tail at a nonzero difference
$b-N$ is bounded by fixed mixed divisor moments
\[
 \sum_{b\asymp x}\tau(b)^R\tau(|b-N|)^S
 \ll x(\log x)^C.
\]
Apply Lemma~\ref{coef:mixed-moment}; its fixed order is included in
$K_{\mathrm{mom}}$.  The exact
difference zero is estimated separately by $N^{\theta_*+o(1)}$.
Choose a sufficiently large logarithmic $H$.  For the retained modes
one pays the Fourier integral norm, rather than the number of modes.

Choose an individual medium slot, if present, in
$X^{3\zeta}\le U\le X^{1/3+\zeta/4}$.
Its local length lies in
$x^{2\zeta}\le U\le x^{1/3+\zeta/2}$ by
\eqref{coef:complete-scales}.  Its sequence has the form
\[
 \alpha_n=e_n\one_I(n)\one_{P^-(n)>z}h(n)n^{it},
 \qquad e_n=1\text{ or }\mu(n),\quad |h|\le1,
\]
with fixed polylogarithmic variation and frequency.  The auxiliary
Siegel--Walfisz property required in the medium estimate is
\begin{equation}
 \left|\sum_{\substack{n\equiv b\ (q)\\(n,k)=1}}\alpha_n
 -\frac1{\varphi(q)}\sum_{(n,kq)=1}\alpha_n\right|
 \ll_A U\tau(k)(\log x)^{-A}.
 \label{coef:rough-SW}
\end{equation}
For completeness, roughness can be inserted at even Brun order
$h=2\lfloor(\log\log x)^2\rfloor$.  Its next divisor support is
$z^{h+1}<U^{1/4}$ and its gap density has arbitrary logarithmic decay.
For squarefree $\ell$,
$\mu(\ell v)=\mu(\ell)\mu(v)\one_{(v,\ell)=1}$.
Put $V_\ell=U/\ell$ for the quotient scale. The identity
\[
 \mu(v)\one_{(v,\ell)=1}
 =\sum_{\substack{w\mid v\\p\mid w\Rightarrow p\mid\ell}}\mu(v/w)
\]
reduces the small-modulus calculation to the unrough
Siegel--Walfisz estimate for $\mu$.  Truncation at $w\le V_\ell^{1/4}$
leaves $O(V_\ell^{7/8}\tau(\ell)^2)$; the smallest inner scale is $U^{9/16}$.
After summing $\ell\le U^{1/4}$ this is
$O(U^{29/32}\log^4U)$.  Large moduli are bounded by the one-bounded
interval estimate.  Inserting $(n,k)=1$ deletes only primes $p\mid k$
above $z$ and costs at most
$2(U/z+1)\omega_{\mathrm{pf}}(k)$.  Partial summation absorbs the allowed variation
and twists by choosing a stronger initial logarithmic accuracy.  This proves \eqref{coef:rough-SW} for the actual slots.

The complementary convolution is bounded by a fixed divisor power.
If necessary, truncate only this complementary coefficient at a fixed
logarithmic threshold, using the same mixed divisor moment bound for
nonzero differences and $N^{\theta_*+o(1)}$ at zero.  The threshold is
chosen after the fixed orders, and the stronger local logarithmic
estimate restores its normalization.  The medium slot is not truncated.
Now apply Proposition~\ref{prop:current-medium} to the exact tensor
representation \eqref{coef:tensor-identity}.

If no medium slot occurs, every original group contains a long
slot: since $\zeta<a_*/100$, six slots below $X^{3\zeta}$ cannot
have product comparable to $R_i$. There are at most two long slots,
because three would exceed the complete scale. A truncated M\"obius
slot is $O(X^{1/3})$ and cannot be long. Thus one or two long
unweighted integer variables remain, with their logarithms, smooth
twists, individual intervals, and common roughness projector.
Their complementary short products are bounded by $x^{20\zeta}$
and $x^{40\zeta}$, respectively. Appendix~\ref{app:brun} estimates
both cases with the exact coprime mean.

Finally, the low-rank coefficient pieces have complete support with a
fixed gap below $x^{1/2}$.  For one original prime use maximal prime
Bombieri--Vinogradov.  For two or more original primes, apply
Lemma~\ref{coef:low-rank-BV} using one original prime slot as explained
above.  The bounded total coefficient norm permits its absolute
form.  All fixed atom scalars, logarithmically many charts and the
Fourier integral norm are paid by requesting stronger local logarithmic
accuracy.  The removed prime powers, collisions and trim have fixed
power savings.  The finite family therefore has one eventual threshold
for the asserted estimate.
\end{proof}

\section{One and two long rough variables}
\label{app:brun}

The estimates here supply the branches without a medium slot in the proof of Proposition~\ref{coef:prime-product-insertion}.

For the one- and two-long-variable cases we use direct progression
estimates and the classical complete Kloosterman bound.
The argument keeps the exact coprime mean and is uniform in the
size of the primitive residue.

\subsection{Brun polynomials with a uniform density gap}

Let $L=\log\log x$, $h=2\lfloor L^2\rfloor$, $r=h+1$ and
$\log z\le2\log x/L^3$.  Put $P(z)=\prod_{p\le z}p$ and
$y=z^r=x^{o(1)}$.  Define
\[
 \lambda^+(d)=\mu(d)\one_{d\mid P(z),\,\omega_{\mathrm{pf}}(d)\le h},
 \qquad
 \lambda^-(d)=\mu(d)\one_{d\mid P(z),\,\omega_{\mathrm{pf}}(d)\le h+1},
\]
and
\[
 U_z(n)=\sum_{d\mid n}\lambda^+(d),\quad
 L_z(n)=\sum_{d\mid n}\lambda^-(d),\quad
 G_z=U_z-L_z,\quad r_z(n)=\one_{(n,P(z))=1}.
\]
Both divisor sums have coefficients of absolute value at most one and
support below $y$.  The alternating binomial identity gives
\begin{equation}
 L_z\le r_z\le U_z,\qquad U_z\ge0,\qquad
 G_z(n)=\binom{\omega_{\mathrm{pf}}((n,P(z)))}{r}\ge0.
 \label{coef:Brun-order}
\end{equation}
This also covers $(n,P(z))=1$, when both polynomials are one.
The lower polynomial need not be nonnegative.

For every $q$ set
\[
 A_q=\sum_{(d,q)=1}\frac{\lambda^+(d)}d,\quad
 g_q=\sum_{(d,q)=1}\frac{\lambda^+(d)-\lambda^-(d)}d,\quad
 V_q=\prod_{\substack{p\le z\\p\nmid q}}\left(1-\frac1p\right).
\]
Finite inclusion-exclusion, applied to independent Bernoulli variables
of means $1/p$, gives
\begin{equation}
 0\le A_q-V_q\le g_q
 \le\frac{(\sum_{p\le z}1/p)^r}{r!}=:\epsilon_x.
 \label{coef:Brun-gap}
\end{equation}
The bound is uniform in the set of primes dividing $q$.  The elementary
prime harmonic bound $\sum_{p\le z}1/p\ll L$ and $r!\ge(r/e)^r$ yield
$\epsilon_x\le\exp(-cL^2\log L)$ eventually.  Thus $\epsilon_x$
is smaller than every fixed negative power of $\log x$, whereas
$y=x^{o(1)}$ will occur only beside a fixed power saving.

\subsection{Smooth progression cores}

For a compactly supported smooth profile define
$N_1(f)=\|f\|_1+\|f'\|_1$ and
$N_2(f)=\|f\|_1+\|f''\|_1$.
Progression summation and M\"obius inversion give
\begin{align}
 \sum_{n\equiv b\ (q)}f(n/U)&=\frac Uq\int f+O(N_1(f)),\notag\\
 \sum_{(n,q)=1}f(n/U)&=U\frac{\varphi(q)}q\int f
                        +O(\tau(q)N_1(f)).
 \label{coef:one-core}
\end{align}
These estimates are uniform in $U>0$ and in the primitive class $b$.
For a fixed $m$ coprime to $q$, the discrepancy of $mn\equiv a\pmod q$
from its exact coprime mean is therefore
$O((1+\tau(q)/\varphi(q))N_1(f))$.  If $(m,q)>1$ and $(a,q)=1$,
both terms vanish.

For two variables, use $e(t)=e^{2\pi it}$ and
$\widehat f(t)=\int f(s)e(-ts)\,ds$.
Let $c_q$ be the Ramanujan sum and
$S(j,k;q)=\sum_{u\bmod q}^{*}e((ju+k\bar u)/q)$.
For $b=a\bar m\pmod q$, double periodization gives the exact
discrepancy
\begin{equation}
 \frac{UV}{q^2}\sum_{j,k\in\mathbb Z}
 \widehat f(jU/q)\widehat g(kV/q)
 \left(S(bj,k;q)-\frac{c_q(j)c_q(k)}{\varphi(q)}\right).
 \label{coef:two-Poisson}
\end{equation}
Both axes vanish identically, because
$S(0,k;q)=c_q(k)$, $S(bj,0;q)=c_q(j)$ and $c_q(0)=\varphi(q)$.
For $jk\ne0$, use
\[
 |S(bj,k;q)|\le\tau(q)\sqrt q\,(j,k,q)^{1/2},\qquad
 |\widehat f(t)|\ll N_2(f)(1+|t|)^{-2}.
\]
For every $H>0$ and $d\ge1$,
$\sum_{j\ne0,\,d\mid j}(1+|j|/H)^{-2}\le2H/d$.
Expanding common divisors and using
$\sum_{d\mid q}d^{-3/2}\le\zeta(3/2)$ therefore bounds
\eqref{coef:two-Poisson} by
\begin{equation}
 N_2(f)N_2(g)
 \left(\sqrt q\,\tau(q)+\frac{\tau(q)^2}{\varphi(q)}\right).
 \label{coef:two-core-bound}
\end{equation}
For the Ramanujan term one uses
$|c_q(j)|\le(j,q)$ and
$\sum_{d\mid q}\varphi(d)/d\le\tau(q)$.
When $q=1$ the discrepancy is zero.

Let a short coefficient $\alpha_m$ have
$\sum|\alpha_m|\ll M\log^C x$, with $x\asymp MU$ or $MUV$.
If $q\le D$ and $|\rho(q)|\le\tau(q)^J\log^B x$ with fixed $J,B$,
the sums of the one- and two-variable core errors are
\begin{equation}
 O(MD\log^{C'}x),\qquad O(MD^{3/2}\log^{C'}x),
 \label{coef:core-summed}
\end{equation}
respectively.  On inserting Brun divisors $d_1,d_2\le y$, the profiles
are at scales $U/d_1,V/d_2$ and their norms in
\eqref{coef:two-core-bound} are unchanged.  Replace $m$ by $md_1d_2$
only in the primitive residue.  The fixed-product bound is uniform in
that residue, so summing the divisor pairs costs at most $y^2$.

\subsection{Restoring roughness with the exact mean}

For a nonnegative smooth $F$, expansion of the divisor coefficients in
\eqref{coef:one-core} gives
\begin{align}
 \sum_{(n,q)=1}F(n/U)U_z(n)
 &=U\frac{\varphi(q)}q A_q\int F+O(y\tau(q)N_1(F)),\notag\\
 \sum_{(n,q)=1}F(n/U)G_z(n)
 &=U\frac{\varphi(q)}q g_q\int F+O(y\tau(q)N_1(F)).
 \label{coef:exact-Brun-means}
\end{align}
Both left sides are nonnegative.  They are exact coprime sums;
the factor $\varphi(q)/q$ is retained throughout.

For two variables write $U_i=U_z(n_i)$, $r_i=r_z(n_i)$ and
$G_i=G_z(n_i)$.  The valid product comparison is
\begin{equation}
 0\le U_1U_2-r_1r_2\le G_1U_2+U_1G_2.
 \label{coef:rough-product-gap}
\end{equation}
Indeed $U_1U_2-r_1r_2=(U_1-r_1)U_2+r_1(U_2-r_2)$.
This argument does not multiply two possibly negative lower polynomials.
For complex profiles choose nonnegative smooth majorants $F\ge|f|$,
$G\ge|g|$ and use $|\alpha|$ only in estimating the absolute gap.
If $C_q$ denotes the congruence count and $M_q$ its exact coprime mean,
and $B=G_1U_2+U_1G_2$, then
\begin{align*}
 |(C_q-M_q)(\alpha fg(r_1r_2-U_1U_2))|
 &\le C_q(|\alpha|FGB)+M_q(|\alpha|FGB)\\
 &\le |(C_q-M_q)(|\alpha|FGB)|+2M_q(|\alpha|FGB).
\end{align*}
The first term is covered by expanding the Brun divisors and
\eqref{coef:core-summed}; the second is estimated by
\eqref{coef:exact-Brun-means}.  Fixed divisor/totient mean bounds give
the total two-variable error
\begin{equation}
 \ll\log^{C'}x\left\{
 MUV\epsilon_x+M(U+V)y+My^2+My^2D^{3/2}\right\}.
 \label{coef:two-rough-error}
\end{equation}
The one-variable version is
\begin{equation}
 \ll\log^{C'}x\{MU\epsilon_x+My+MyD\}.
 \label{coef:one-rough-error}
\end{equation}
Only fixed logarithmic factors multiply $\epsilon_x$.
For the long alternatives, $M\le x^{20\zeta}$ or $x^{40\zeta}$,
$D\le x^{\Theta_{\mathrm{loc}}}$ with $\Theta_{\mathrm{loc}}<3/5$,
and $\min(U,V)\ge x^\nu$ for a fixed $\nu>0$.
Choose $\zeta$ so that
\[
 20\zeta+\Theta_{\mathrm{loc}}<1,\qquad
 40\zeta+\frac32\Theta_{\mathrm{loc}}<1.
\]
Every remaining term in \eqref{coef:two-rough-error} and
\eqref{coef:one-rough-error} then has a fixed power saving.
The errors are consequently $O_A(x\log^{-A}x)$ for every fixed $A$.

\subsection{Individual sharp intervals}

Let $\delta=(\log x)^{-B}$.  An interval in a fixed positive scale has
a smooth approximation whose error is bounded by nonnegative smooth
strips of normalized width $O(\delta)$ at its endpoints.  If the interval
is shorter than $\delta$, a single strip covers it.  Their integrals
are $O(\delta)$, their first-derivative norms are $O(1)$, and their
second-derivative norms are bounded by a fixed power of $\delta^{-1}$.
Multiply by the fixed envelopes for the other profiles.

Apply the proved rough smooth estimates to these positive strips.
Their exact means are bounded by
\[
 \ll\log^{C'}x\{x\delta+M(U+V)y+My^2\}
\]
in the two-variable case, and by
$\log^{C'}x\{x\delta+My\}$ in the one-variable case.
The exponent in the first term uses the strip integral, not its
derivatives, and is independent of $B$.  The $B$-dependent derivative
losses affect only the fixed-power errors.  Choose $B$ after the desired
logarithmic accuracy.  This controls both the count and the exact mean,
uniformly in every primitive residue, including a residue equal to an
intermediate complete product.  It proves the required extension to
sharp individual intervals.  No assertion about arbitrary sharp coupled
domains is used.
